%   terminologia: 
%     processo: refere-se à execução no espaço do usuário, ou talvez à struct proc etc.
%     memória do processo: a parte inferior do espaço de endereçamento
%     o processo possui uma thread com duas pilhas (uma para modo kernel e uma para 
%     modo usuário)
% breve explicação sobre as condições iniciais da tabela de páginas:
%     paginação desligada, mas o mapeamento virtual é em sua maioria direto para o físico,
%     que é como as coisas parecem quando ligamos a paginação também,
%     já que a paginação é ativada após criarmos o primeiro processo.
%   mencionar por que ainda existem SEG_UCODE/SEG_UDATA?
%   será que realmente explicamos o que os dois bits inferiores de %cs fazem?
%     em particular sua interação com PTE_U
%   nota lateral sobre por que é extern char[]
\chapter{Organização do sistema operacional} \label{CH:FIRST}
Um requisito fundamental para um sistema operacional é suportar várias atividades simultaneamente. Por exemplo, usando a interface de chamadas de sistema descrita no Capítulo~\ref{CH:UNIX}, um processo pode iniciar novos processos com \lstinline{fork}. O sistema operacional deve \indextext{compartilhar o tempo} dos recursos do computador entre esses processos. Por exemplo, mesmo que haja mais processos do que CPUs físicas, o sistema operacional deve garantir que todos os processos tenham a chance de executar. O sistema operacional também deve providenciar \indextext{isolamento} entre os processos. Ou seja, se um processo tiver um bug e apresentar mau funcionamento, ele não deve afetar processos que não dependem do processo com erro. Contudo, isolamento completo é excessivo, já que deve ser possível que processos interajam intencionalmente; pipelines são um exemplo disso. Assim, um sistema operacional deve cumprir três requisitos: multiplexação, isolamento e interação.
Este capítulo fornece uma visão geral de como sistemas operacionais são organizados para atingir esses três requisitos. Existem várias maneiras de fazer isso, mas este texto foca em projetos tradicionais centrados em um \indextext{kernel monolítico}, utilizado por muitos sistemas Unix. Este capítulo também oferece uma visão geral de um processo no xv6, que é a unidade de isolamento no xv6, e a criação do primeiro processo quando o xv6 inicia.
O xv6 roda em um microprocessador RISC-V \indextext{multi-core}\footnote{ Por "multi-core", este texto se refere a múltiplas CPUs que compartilham memória, mas executam em paralelo, cada uma com seu próprio conjunto de registradores. Este texto às vezes usa o termo \indextext{multiprocessador} como sinônimo de multi-core, embora multiprocessador também possa se referir mais especificamente a um computador com vários chips de processador distintos.} e grande parte de sua funcionalidade de baixo nível (por exemplo, a implementação de processos) é específica do RISC-V. O RISC-V é uma CPU de 64 bits, e o xv6 é escrito em C "LP64", o que significa que long (L) e ponteiros (P) na linguagem C têm 64 bits, mas um {\tt int} tem 32 bits. Este livro assume que o leitor já tenha feito alguma programação em nível de máquina em alguma arquitetura, e introduzirá as ideias específicas do RISC-V à medida que forem surgindo. Os documentos da ISA de nível de usuário~\cite{riscv:user} e da arquitetura privilegiada~\cite{riscv:priv} são as especificações completas. Você também pode consultar "The RISC-V Reader: An Open Architecture Atlas"~\cite{riscv}.
A CPU em um computador completo é cercada por hardware de suporte, em grande parte na forma de interfaces de E/S. O xv6 é escrito para o hardware de suporte simulado pela opção "-machine virt" do qemu. Isso inclui RAM, uma ROM contendo código de boot, uma conexão serial com o teclado/tela do usuário e um disco para armazenamento.

%%
\section{Abstraindo recursos físicos} 
%%

A primeira pergunta que alguém pode fazer ao se deparar com um sistema operacional é: por que tê-lo? Ou seja, poderíamos implementar as chamadas de sistema da Figura~\ref{fig:api} como uma biblioteca, com a qual os aplicativos são vinculados. Nesse plano, cada aplicação poderia até ter sua própria biblioteca adaptada às suas necessidades. As aplicações poderiam interagir diretamente com os recursos de hardware e usá-los da melhor forma possível para alcançar alto desempenho ou desempenho previsível. Alguns sistemas operacionais para dispositivos embarcados ou sistemas em tempo real são organizados dessa forma.
O problema dessa abordagem com bibliotecas é que, se houver mais de uma aplicação rodando, elas devem se comportar bem. Por exemplo, cada aplicação deve ceder periodicamente a CPU para que outras possam executar. Esse esquema de compartilhamento de tempo \textit{cooperativo} pode funcionar se todas as aplicações confiarem umas nas outras e não tiverem bugs. Mas o mais comum é que as aplicações não confiem umas nas outras e tenham bugs, então geralmente se deseja um isolamento mais forte do que o fornecido por esse esquema cooperativo.
Para alcançar um isolamento forte, é útil proibir as aplicações de acessarem diretamente recursos sensíveis de hardware, abstraindo esses recursos em serviços. Por exemplo, aplicações Unix interagem com armazenamento apenas por meio das chamadas de sistema do sistema de arquivos, como \lstinline{open}, \lstinline{read}, \lstinline{write} e \lstinline{close}, ao invés de ler e escrever diretamente no disco. Isso fornece à aplicação a conveniência de nomes de caminhos, e permite que o sistema operacional (como implementador da interface) gerencie o disco. Mesmo que o isolamento não seja uma preocupação, programas que interagem intencionalmente (ou que apenas queiram se manter separados) provavelmente acharão o sistema de arquivos uma abstração mais conveniente do que o uso direto do disco.
Da mesma forma, o Unix troca CPUs de hardware entre os processos de forma transparente, salvando e restaurando o estado dos registradores conforme necessário, de modo que as aplicações não precisam estar cientes do compartilhamento de tempo. Essa transparência permite que o sistema operacional compartilhe as CPUs mesmo que algumas aplicações estejam em loops infinitos.
Outro exemplo: processos Unix usam \lstinline{exec} para montar sua imagem de memória, ao invés de interagir diretamente com a memória física. Isso permite que o sistema operacional decida onde posicionar um processo na memória; se a memória estiver escassa, o sistema operacional pode até armazenar parte dos dados do processo em disco. \lstinline{exec} também oferece aos usuários a conveniência de um sistema de arquivos para armazenar imagens executáveis de programas.
Muitas formas de interação entre processos Unix ocorrem por meio de descritores de arquivos. Esses descritores não apenas abstraem muitos detalhes (por exemplo, onde os dados em um pipe ou arquivo estão armazenados), mas também são definidos de forma a simplificar a interação. Por exemplo, se uma aplicação em um pipeline falhar, o kernel gera um sinal de fim de arquivo para o próximo processo no pipeline.
A interface de chamadas de sistema da Figura~\ref{fig:api} foi cuidadosamente projetada para fornecer conveniência ao programador e a possibilidade de isolamento forte. A interface Unix não é a única maneira de abstrair recursos, mas tem se mostrado uma abordagem eficaz.

%% 
\section{Modo usuário, modo supervisor e chamadas de sistema}
%% 

Um isolamento forte requer uma barreira rígida entre as aplicações e o sistema
operacional. Se a aplicação comete um erro, não queremos que o sistema
operacional falhe, nem que outras aplicações sejam afetadas. Em vez disso, o sistema
operacional deve ser capaz de encerrar a aplicação com falha e continuar executando
as demais. Para alcançar esse isolamento, o sistema operacional deve garantir que
as aplicações não possam modificar (ou sequer ler) as estruturas de dados e
instruções do sistema operacional, nem acessar a memória de outros processos.

As CPUs fornecem suporte de hardware para esse tipo de isolamento. Por exemplo,
o RISC-V possui três modos em que a CPU pode executar instruções:
\indextext{modo máquina},
\indextext{modo supervisor} e
\indextext{modo usuário}.
Instruções executadas em modo máquina têm privilégios totais; a
CPU inicia em modo máquina. Esse modo é usado principalmente para
configurar o computador durante o boot. O xv6 executa algumas poucas linhas em
modo máquina e depois muda para o modo supervisor.

No modo supervisor, a CPU tem permissão para executar 
\indextext{instruções privilegiadas}:
por exemplo, habilitar e desabilitar interrupções, ler e escrever
no registrador que contém o endereço da tabela de páginas, etc.
Se uma aplicação no modo usuário tentar executar
uma instrução privilegiada, a CPU não a executa; em vez disso, muda
para o modo supervisor para que o código em modo supervisor possa encerrar a aplicação,
pois ela tentou fazer algo proibido.
A Figura~\ref{fig:os}
no Capítulo~\ref{CH:UNIX} ilustra essa organização. Uma aplicação pode
executar apenas instruções do modo usuário (por exemplo, somar números etc.) e diz-se que
está em execução no 
\indextext{espaço do usuário},
enquanto o software em modo supervisor pode também executar instruções privilegiadas e
diz-se que está executando no
\indextext{espaço do kernel}.
O software que roda no espaço do kernel (ou em modo supervisor) é chamado de
\indextext{kernel}.

Uma aplicação que deseja invocar uma função do kernel (por exemplo, a
chamada de sistema \lstinline{read} no xv6) precisa
fazer uma transição para o kernel; uma aplicação \emph{não pode} invocar
uma função do kernel diretamente. As CPUs oferecem uma instrução especial que troca
o modo da CPU de usuário para supervisor e entra no kernel por um ponto de entrada
especificado pelo próprio kernel. (O RISC-V
fornece a instrução
\indexcode{ecall}
para esse fim.) Uma vez que a CPU mudou para o modo supervisor,
o kernel pode então validar os argumentos da chamada de sistema (por exemplo,
verificar se o endereço passado está dentro da memória da aplicação), decidir se
a aplicação tem permissão para realizar a operação solicitada (por exemplo,
verificar se ela pode escrever no arquivo especificado), e então negar
ou executar a operação. É fundamental que o kernel controle o ponto de entrada para
as transições ao modo supervisor; se a aplicação pudesse decidir esse ponto de entrada,
uma aplicação maliciosa poderia, por exemplo, entrar no kernel em um ponto onde a
validação dos argumentos fosse ignorada.

%% 
\section{Organização do kernel}
%% 

Uma questão fundamental de projeto é: que parte do sistema
operacional deve ser executada em modo supervisor?
Uma possibilidade é que todo o sistema operacional resida
no kernel, de modo que as implementações de todas as chamadas de sistema
executem em modo supervisor.
Essa organização é chamada de
\indextext{kernel monolítico}.

Nessa organização, todo o sistema operacional consiste em um único
programa executando com privilégios completos de hardware.
Essa abordagem é conveniente porque o projetista do sistema operacional não precisa
decidir quais partes do sistema não precisam de acesso total ao hardware.
Além disso, é mais fácil para as diferentes partes do sistema operacional
cooperarem. Por exemplo, o sistema pode ter um cache de blocos que seja
compartilhado tanto pelo sistema de arquivos quanto pelo sistema de memória virtual.

Uma desvantagem da organização monolítica é que as interações entre
as diferentes partes do sistema operacional costumam ser complexas (como veremos
ao longo deste texto), e, por isso, é fácil que o desenvolvedor cometa erros.
Em um kernel monolítico, um erro pode ser fatal, pois uma falha em modo supervisor
geralmente causa a falha do kernel. Se o kernel falhar,
o computador para de funcionar e todas as aplicações falham também.
É necessário reiniciar o computador para recomeçar.

Para reduzir o risco de erros no kernel, projetistas de sistemas operacionais podem minimizar a
quantidade de código do sistema que roda em modo supervisor, executando a
maior parte do sistema operacional em modo usuário.
Essa organização do kernel é chamada de
\indextext{microkernel}.

\begin{figure}[t]
\center
\includegraphics[scale=0.5]{fig/mkernel.pdf}
\caption{Um microkernel com um servidor de sistema de arquivos}
\label{fig:mkernel}
\end{figure}

A Figura~\ref{fig:mkernel}
ilustra esse projeto de microkernel. Na figura, o sistema de arquivos é executado como
um processo em nível de usuário. Os serviços do sistema operacional executados como processos são chamados de servidores.
Para permitir que aplicações interajam com o
servidor de arquivos, o kernel fornece um mecanismo de comunicação entre processos
para enviar mensagens de um processo em modo usuário para outro. Por exemplo, se uma aplicação como o shell
deseja ler ou escrever um arquivo, ela envia uma mensagem ao servidor de arquivos e espera
uma resposta.

Em um microkernel, a interface do kernel consiste em algumas funções de baixo nível
para iniciar aplicações, enviar mensagens,
acessar hardware de dispositivos, etc. Essa organização permite que o kernel seja
relativamente simples, já que a maior parte do sistema operacional
reside em servidores em modo usuário.

Na prática, tanto kernels monolíticos quanto microkernels são
populares. Muitos
kernels Unix são monolíticos. Por exemplo, o Linux possui um kernel monolítico,
embora algumas funções do sistema operacional sejam executadas como servidores em modo usuário
(por exemplo, o sistema de janelas). O Linux oferece alto desempenho para aplicações intensivas em SO,
em parte porque os subsistemas do kernel podem ser fortemente integrados.

Sistemas operacionais como Minix, L4 e QNX são organizados como microkernels com
servidores, e são amplamente utilizados em ambientes embarcados.
Uma variante do L4, o seL4, é pequeno o suficiente para ter sido verificado formalmente quanto à
segurança de memória e outras propriedades de segurança~\cite{sel4}.

Há muito debate entre os desenvolvedores de sistemas operacionais sobre qual
organização é melhor, e não há evidência conclusiva a favor de uma ou outra.
Além disso, tudo depende do que se entende por “melhor”:
desempenho mais rápido, tamanho menor do código, confiabilidade do kernel,
confiabilidade do sistema como um todo (incluindo os serviços em modo usuário), etc.

Também existem considerações práticas que podem ser mais importantes
do que a questão da organização em si. Alguns sistemas operacionais
têm microkernel, mas executam alguns dos serviços em modo usuário no espaço do kernel
por razões de desempenho. Alguns sistemas têm kernels monolíticos
porque começaram assim, e não há muito incentivo para migrar para uma organização
puramente baseada em microkernel, já que novos recursos podem ser mais importantes
do que reescrever o sistema existente para se adequar a um design de microkernel.

Do ponto de vista deste livro, sistemas operacionais microkernel e monolíticos
compartilham muitas ideias fundamentais. Ambos implementam chamadas de sistema, usam
tabelas de páginas, tratam interrupções, suportam processos, utilizam
locks para controle de concorrência, implementam um sistema de arquivos,
etc. Este livro foca nessas ideias centrais.

O xv6 é
implementado como um kernel monolítico, como a maioria dos sistemas Unix.
Assim, a interface do kernel do xv6 corresponde à interface do sistema operacional,
e o kernel implementa o sistema completo. Como
o xv6 não fornece muitos serviços, seu kernel é menor do que o de alguns
microkernels, mas conceitualmente o xv6 é monolítico.

\section{Código: organização do xv6}

\begin{figure}[t]
\center
\begin{tabular}{l|l}
{\bf Arquivo} & {\bf Descrição} \\
\midrule
bio.c & Cache de blocos de disco para o sistema de arquivos. \\
console.c & Conexão com o teclado e a tela do usuário. \\
entry.S & Primeiras instruções de boot. \\
exec.c & Implementação da chamada de sistema exec(). \\
file.c & Suporte a descritores de arquivos. \\
fs.c & Sistema de arquivos. \\
kalloc.c & Alocador de páginas físicas. \\
kernelvec.S & Tratamento de traps no kernel. \\
log.c & Log de sistema de arquivos e recuperação de falhas. \\
main.c & Inicialização dos módulos durante o boot. \\
pipe.c & Pipes. \\
plic.c & Controlador de interrupções RISC-V. \\
printf.c & Saída formatada para o console. \\
proc.c & Processos e escalonamento. \\
sleeplock.c & Locks que liberam a CPU. \\
spinlock.c & Locks que não liberam a CPU. \\
start.c & Código inicial em modo máquina. \\
string.c & Biblioteca de strings e arrays de bytes em C. \\
swtch.S & Troca de contexto de threads. \\
syscall.c & Encaminha chamadas de sistema para as funções correspondentes. \\
sysfile.c & Chamadas de sistema relacionadas a arquivos. \\
sysproc.c & Chamadas de sistema relacionadas a processos. \\
trampoline.S & Código em assembly para alternar entre usuário e kernel. \\
trap.c & Código C para tratar e retornar de traps e interrupções. \\
uart.c & Driver de console da porta serial. \\
virtio\_disk.c & Driver de dispositivo de disco. \\
vm.c & Gerenciamento de tabelas de páginas e espaços de endereçamento. \\
\end{tabular}
\caption{Arquivos fonte do kernel xv6.}
\label{fig:source}
\end{figure}

O código-fonte do kernel do xv6 está no subdiretório {\tt kernel/}.
O código é dividido em arquivos, seguindo uma ideia geral de modularidade;
a Figura~\ref{fig:source} lista os arquivos. As interfaces entre os módulos
são definidas em \lstinline{defs.h} \fileref{kernel/defs.h}.

%% 
\section{Visão geral do processo}
%% 

A unidade de isolamento no xv6 (assim como em outros sistemas operacionais Unix) é um 
\indextext{processo}.
A abstração de processo impede que um processo destrua ou espione
a memória, CPU, descritores de arquivos, etc., de outro processo. Também impede que um processo
corrompa o próprio kernel, de modo que ele não possa subverter os mecanismos de isolamento do kernel.
O kernel deve implementar a abstração de processo com cuidado, pois
um aplicativo com bugs ou malicioso pode tentar enganar o kernel ou o hardware para fazer
algo prejudicial (por exemplo, burlar o isolamento). Os mecanismos usados pelo
kernel para implementar processos incluem a flag de modo usuário/supervisor, os espaços de endereçamento
e a divisão de tempo (time-slicing) entre threads.

Para ajudar a reforçar o isolamento, a abstração de processo fornece a
ilusão ao programa de que ele possui sua própria máquina privada. Um processo fornece
ao programa o que parece ser um sistema de memória próprio, ou
\indextext{espaço de endereçamento},
ao qual outros processos não têm acesso de leitura ou escrita.
Um processo também fornece ao programa o que parece ser sua própria
CPU para executar as instruções do programa.

O xv6 utiliza tabelas de páginas (implementadas pelo hardware) para fornecer a cada processo
seu próprio espaço de endereçamento. A tabela de páginas do RISC-V
traduza (ou ``mapeia'') um
\indextext{endereço virtual}
(o endereço que uma instrução RISC-V manipula) para um
\indextext{endereço físico}
(um endereço que a CPU envia para a memória principal).

\begin{figure}[t]
\centering
\includegraphics[scale=0.5]{fig/as.pdf}
\caption{Organização do espaço de endereçamento virtual de um processo}
\label{fig:as}
\end{figure}

O xv6 mantém uma tabela de páginas separada para cada processo que define o espaço de
endereçamento desse processo. Como ilustrado na
Figura~\ref{fig:as},
um espaço de endereçamento inclui a
\indextext{memória de usuário}
do processo começando no endereço virtual zero. As instruções vêm primeiro,
seguidas por variáveis globais, depois a pilha (stack),
e finalmente uma área de "heap" (para malloc)
que o processo pode expandir conforme necessário.
Existem vários fatores que limitam o
tamanho máximo do espaço de endereçamento de um processo:
os ponteiros no RISC-V têm 64 bits de largura;
o hardware usa apenas os 39 bits inferiores ao
consultar endereços virtuais nas tabelas de páginas;
e o xv6 usa apenas 38 desses 39 bits.
Assim, o endereço máximo é $2^{38}-1$ =
0x3fffffffff, que é \lstinline{MAXVA}~\lineref{kernel/riscv.h:/define.MAXVA/}.
No topo do espaço de endereçamento, o xv6 coloca
uma página de \indextext{trampolim} (4096 bytes)
e uma página de \indextext{trapframe}. O xv6 utiliza essas duas páginas
para fazer a transição para o kernel e de volta;
a página de trampolim contém o código para fazer a transição de entrada e saída
do kernel, e o trapframe é onde o kernel salva
os registradores de usuário do processo,
como será explicado no Capítulo~\ref{CH:TRAP}.

O kernel do xv6 mantém muitas informações de estado para cada processo,
que são agrupadas em uma
\indexcode{struct proc}
\lineref{kernel/proc.h:/^struct.proc/}.
As informações mais importantes de estado de kernel de um processo são sua
tabela de páginas, sua pilha de kernel e seu estado de execução.
Usaremos a notação
\indexcode{p->xxx}
para nos referirmos aos elementos da estrutura
\lstinline{proc};
por exemplo,
\indexcode{p->pagetable} é um ponteiro para a tabela de páginas do processo.

Cada processo possui um fluxo de controle (ou
\indextext{thread}, em inglês) que mantém o estado necessário para executar o processo.
Em determinado momento, uma thread pode estar executando em uma CPU,
ou suspensa (não executando, mas capaz de retomar a execução
no futuro).
Para alternar uma CPU entre processos,
o kernel suspende a thread que está atualmente em execução nessa CPU
e salva seu estado,
e restaura o estado da thread anteriormente suspensa de outro processo.
Grande parte do estado de uma thread (variáveis locais, endereços de retorno de chamadas de função)
é armazenado nas pilhas da thread.
Cada processo possui duas pilhas: uma pilha de usuário e uma pilha de kernel
(\indexcode{p->kstack}).
Quando o processo está executando instruções de usuário, apenas sua pilha de usuário
está em uso, e sua pilha de kernel está vazia.
Quando o processo entra no kernel (por uma chamada de sistema ou interrupção),
o código do kernel é executado na pilha de kernel do processo; enquanto
o processo está no kernel, sua pilha de usuário ainda contém dados salvos,
mas não está sendo usada ativamente.
A thread de um processo alterna entre o uso ativo da pilha de usuário
e da pilha de kernel. A pilha de kernel é separada (e protegida do
código de usuário) para que o kernel
possa executar mesmo que o processo tenha corrompido sua pilha de usuário.

Um processo pode fazer uma chamada de sistema executando a instrução \indexcode{ecall}
do RISC-V. Essa instrução eleva o nível de privilégio do hardware e
altera o contador de programa para um ponto de entrada definido pelo kernel. O código
nesse ponto de entrada alterna para a pilha de kernel do processo e executa as instruções do kernel
que implementam a chamada de sistema. Quando a chamada de sistema
é concluída, o kernel alterna de volta para a pilha de usuário e retorna ao
modo usuário chamando a instrução \indexcode{sret}, que reduz
o nível de privilégio do hardware e retoma a execução das instruções de usuário
logo após a instrução da chamada de sistema. A thread de um processo pode
"bloquear" no kernel enquanto espera por I/O, e retomar de onde parou
quando o I/O estiver concluído.

\indexcode{p->state}
indica se o processo está alocado, pronto
para executar, atualmente executando em uma CPU, aguardando por I/O, ou finalizando.

\indexcode{p->pagetable}
armazena a tabela de páginas do processo, no formato
esperado pelo hardware RISC-V.
O xv6 faz com que o hardware de paginação utilize a
\lstinline{p->pagetable}
ao executar aquele processo no modo usuário.
A tabela de páginas de um processo também serve como o registro dos
endereços das páginas físicas alocadas para armazenar a memória do processo.

Resumindo, um processo combina duas ideias de projeto: um espaço de endereçamento para
dar ao processo a ilusão de possuir sua própria memória, e uma thread para dar
ao processo a ilusão de possuir sua própria CPU. No xv6, um processo consiste
em um espaço de endereçamento e uma thread. Em sistemas operacionais reais,
um processo pode ter mais de uma thread para aproveitar múltiplas CPUs.

%%
\section{Código: iniciando o xv6, o primeiro processo e chamada de sistema}
%%

Para tornar o xv6 mais concreto, vamos esboçar como o kernel inicia e executa o primeiro processo. Os capítulos seguintes descreverão em mais detalhes os mecanismos apresentados nesta visão geral.

Quando o computador RISC-V é ligado, ele se inicializa e executa um carregador de boot armazenado na memória somente leitura. O carregador de boot carrega o kernel do xv6 na memória. Em seguida, no modo máquina, a CPU executa o xv6 a partir de
\indexcode{_entry}
\lineref{kernel/entry.S:/^.entry:/}.
O RISC-V começa com o hardware de paginação desabilitado: os endereços virtuais mapeiam diretamente para endereços físicos.

O carregador carrega o kernel do xv6 na memória no endereço físico
\texttt{0x80000000}.
A razão para colocá-lo em
\texttt{0x80000000}
em vez de
\texttt{0x0}
é que a faixa de endereços
\texttt{0x0:0x80000000}
contém dispositivos de I/O.

As instruções em
\lstinline{_entry}
configuram uma pilha para que o xv6 possa executar código C.
O xv6 declara espaço para uma pilha inicial,
\lstinline{stack0},
no arquivo
\lstinline{start.c}
\lineref{kernel/start.c:/stack0/}.
O código em
\lstinline{_entry}
carrega o registrador de ponteiro de pilha
\texttt{sp}
com o endereço
\lstinline{stack0+4096},
o topo da pilha, porque a pilha no RISC-V cresce para baixo.
Agora que o kernel tem uma pilha,
\lstinline{_entry}
chama o código C em
\lstinline{start}
\lineref{kernel/start.c:/^start/}.

A função
\lstinline{start}
realiza algumas configurações que só são permitidas no modo máquina, e então muda para o modo supervisor.
Para entrar no modo supervisor, o RISC-V fornece a instrução
\lstinline{mret}.
Essa instrução normalmente é usada para retornar de uma chamada anterior do modo supervisor para o modo máquina.
\lstinline{start} não está retornando de tal chamada, mas sim configurando as coisas como se estivesse:
ela define o modo de privilégio anterior como supervisor no registrador
\lstinline{mstatus},
define o endereço de retorno como
\lstinline{main}
escrevendo o endereço de
\lstinline{main}
no registrador
\lstinline{mepc},
desativa a tradução de endereços virtuais no modo supervisor escrevendo
\lstinline{0}
no registrador de tabela de páginas
\lstinline{satp},
e delega todas as interrupções e exceções ao modo supervisor.

Antes de saltar para o modo supervisor,
\lstinline{start}
realiza uma última tarefa: programa o chip de clock para gerar interrupções de temporizador.
Com essa configuração finalizada,
\lstinline{start}
“retorna” ao modo supervisor chamando
\lstinline{mret}.
Isso faz com que o contador de programa mude para
\lstinline{main}
\lineref{kernel/main.c:/^main/},
o endereço armazenado anteriormente em \lstinline{mepc}.

Após
\lstinline{main}
\lineref{kernel/main.c:/^main/}
inicializar vários dispositivos e subsistemas, ele cria o primeiro processo chamando
\lstinline{userinit}
\lineref{kernel/proc.c:/^userinit/}.
O primeiro processo executa um pequeno programa escrito em assembly RISC-V, que faz a primeira chamada de sistema no xv6.
\indexcode{initcode.S}
\lineref{user/initcode.S:3}
carrega o número da chamada de sistema \lstinline{exec},
\lstinline{SYS_EXEC}
\lineref{kernel/syscall.h:/exec/},
no registrador {\tt a7}, e então chama \lstinline{ecall} para reentrar no kernel.

O kernel usa o número no registrador {\tt a7} em \lstinline{syscall}
\lineref{kernel/syscall.c:/^syscall/}
para chamar a chamada de sistema desejada.
A tabela de chamadas de sistema
\lineref{kernel/syscall.c:/syscalls/}
mapeia \lstinline{SYS_EXEC} para a função \lstinline{sys_exec}, que o kernel invoca.
Como vimos no Capítulo~\ref{CH:UNIX}, \indexcode{exec} substitui a memória e os registradores do processo atual por um novo programa (neste caso, \indexcode{/init}).

Uma vez que o kernel tenha concluído \lstinline{exec},
ele retorna ao espaço de usuário no processo
\lstinline{/init}.
\lstinline{init}
\lineref{user/init.c:/^main/}
cria um novo arquivo de dispositivo de console, se necessário, e então o abre como os descritores de arquivo 0, 1 e 2.
Depois disso, inicia um shell no console.
O sistema está no ar.

%%
\section{Modelo de Segurança}
%%

Você pode se perguntar como o sistema operacional lida com código com bugs ou malicioso. Como lidar com código malicioso é mais difícil do que com bugs acidentais, é razoável focar principalmente em fornecer segurança contra ações maliciosas. Aqui está uma visão geral das suposições e objetivos típicos de segurança no projeto de sistemas operacionais.

O sistema operacional deve assumir que o código de nível de usuário de um processo fará o possível para danificar o kernel ou outros processos. O código de usuário pode tentar desreferenciar ponteiros fora de seu espaço de endereçamento permitido; pode tentar executar qualquer instrução RISC-V, mesmo aquelas não destinadas ao código de usuário; pode tentar ler e escrever em qualquer registrador de controle do RISC-V; pode tentar acessar diretamente o hardware dos dispositivos; e pode passar valores engenhosos para chamadas de sistema tentando enganar o kernel para travar ou fazer algo incorreto. O objetivo do kernel é restringir cada processo de usuário para que ele só possa ler/escrever/executar sua própria memória de usuário, usar os 32 registradores de propósito geral do RISC-V, e afetar o kernel e outros processos apenas nas formas previstas pelas chamadas de sistema. O kernel deve impedir qualquer outra ação. Isso é tipicamente um requisito absoluto no projeto de kernels.

As expectativas para o próprio código do kernel são bastante diferentes. Assume-se que o código do kernel seja escrito por programadores bem-intencionados e cuidadosos. Espera-se que esse código esteja livre de bugs, e certamente que não contenha nada malicioso. Essa suposição afeta como analisamos o código do kernel. Por exemplo, há muitas funções internas do kernel (como os spinlocks) que causariam sérios problemas se usadas incorretamente. Ao examinar qualquer trecho específico de código do kernel, queremos nos convencer de que ele se comporta corretamente. Assumimos, no entanto, que o código do kernel em geral é escrito corretamente e segue todas as regras de uso das funções e estruturas de dados do próprio kernel. Em nível de hardware, assume-se que a CPU RISC-V, RAM, disco, etc. operem conforme descrito na documentação, sem bugs de hardware.

Claro que na vida real as coisas não são tão simples. É difícil impedir que código de usuário engenhoso torne o sistema inutilizável (ou leve a um pânico) ao consumir recursos protegidos pelo kernel — espaço em disco, tempo de CPU, slots na tabela de processos, etc. É geralmente impossível escrever um kernel sem bugs ou projetar hardware sem defeitos; se os autores de código malicioso conhecerem bugs no kernel ou no hardware, irão explorá-los. Mesmo em kernels maduros e amplamente utilizados, como o Linux, vulnerabilidades continuam sendo descobertas regularmente~\cite{mitre:cves}. Vale a pena projetar salvaguardas no kernel contra a possibilidade de ele conter bugs: afirmações (assertions), verificação de tipos, páginas de guarda na pilha, etc. Finalmente, a distinção entre código de usuário e de kernel às vezes é difusa: alguns processos privilegiados em nível de usuário podem fornecer serviços essenciais e funcionar efetivamente como parte do sistema operacional, e em alguns sistemas operacionais o código de usuário privilegiado pode inserir novo código no kernel (como nos módulos carregáveis do kernel no Linux).

%% 
\section{Mundo real}
%% 

A maioria dos sistemas operacionais adotou o conceito de processo, e a maioria
dos processos se parece com os do xv6. No entanto, sistemas operacionais modernos
suportam vários threads dentro de um processo, para permitir que um único processo
aproveite múltiplas CPUs. Suportar múltiplos threads em um
processo envolve uma quantidade considerável de mecanismos que o xv6 não possui,
frequentemente incluindo mudanças na interface (por exemplo, o
\lstinline{clone}
do Linux, uma variação de
\lstinline{fork}),
para controlar quais aspectos
de um processo os threads compartilham.

%% 
\section{Exercícios}
%%

\begin{enumerate}

\item Adicione uma chamada de sistema ao xv6
  que retorne a quantidade de memória livre disponível.

% break *0x3ffffff000
% disas 0x3ffffff000,+8
% set disassemble-next-line auto
% x/i $pc
% break *0x3ffffff10e
% print/x $sepc
% print/x $pc

\end{enumerate}

